\section{Example Bank: How Frameworks Actually Form}

\subsection{Example 1: Study Entry Ritual}

\begin{example}[Repeated entry cue creates cheap re-entry]
A student repeatedly begins work at the same place, with the same notebook
opened to the next unresolved line, and the same first five-minute task. Over
time the study state becomes easier to enter after small interruptions.
\end{example}

\begin{discussion}
This is a clean case of emergence through repeated coupling. The local events
are simple: same location, same cue, same entry action. The higher-level
result is a framework in which study is no longer reconstructed from scratch.
The evidence lane here is \textbf{medium}: complexity and emergence texts plus
the metastability review justify the structure of the explanation, but no
strong empirical paper in the current source set directly proves this exact
study ritual.
\end{discussion}

\subsection{Example 2: Phone Checking As Interference Architecture}

\begin{example}[Competing route gains historical priority]
Suppose a person repeatedly checks the phone at every short idle boundary:
after a paragraph, after a meal, while standing in line, before sleep. Over
time the idle-to-phone route becomes faster, more expected, and more
self-reinforcing than idle-to-work resumption.
\end{example}

\begin{discussion}
This example shows interference rather than simple habit accumulation. The old
route now intrudes into unrelated contexts because it has acquired superior
re-entry economy. In system terms, the phone route has become a highly
accessible neighbouring basin. The explanatory value is \textbf{medium} and
comes from path dependence, emergence, and attractor language rather than from
any claim of direct neural decoding.
\end{discussion}

\subsection{Example 3: Exercise Identity Through Multiscale Organization}

\begin{example}[Framework forms across body, schedule, and social timing]
A worker begins going to the gym on three fixed evenings, lays out clothes in
advance, travels with the same route, and meets the same partner at the same
time. After several weeks, the ``exercise evening'' no longer feels like a
fresh moral decision each time.
\end{example}

\begin{discussion}
This case is important because it is multiscale. The framework is not formed by
inner resolve alone. Body preparation, route repetition, calendar structure,
and social synchronization all contribute. The macro-pattern becomes
explanatory because it changes what the evening \emph{is}. The system now
carries a regime, not merely a wish.
\end{discussion}

\subsection{Example 4: Research Writing And Negative Path Dependence}

\begin{example}[Avoidance becomes organized competence]
A graduate student repeatedly resolves uncertainty by reading one more paper,
making one more note, or polishing citations rather than writing draft prose.
Eventually the pre-writing state becomes highly skilled, while visible writing
feels increasingly costly.
\end{example}

\begin{discussion}
This is a good example of negative emergence. The system does not merely fail
to write; it becomes organized around avoiding exposure. The resulting
framework has explanatory value because it changes later cost structure. The
manuscript should therefore not ask only ``Why is the student reluctant
today?'' but also ``What repeated micro-policy has trained this reluctance into
a stable route?''
\end{discussion}

\subsection{Example 5: Group Norm Formation}

\begin{example}[Frameworks can be collective]
In a team meeting, people repeatedly reward quick agreement and implicitly
punish slow clarification. After enough repetitions, the group develops a
framework in which fast consensus feels natural and careful dissent feels
socially expensive.
\end{example}

\begin{discussion}
This example broadens the book's relevance. Framework formation is not only
personal. Repeated local interactions can generate collective operating
patterns that later govern what individuals experience as easy or costly.
Complexity texts and emergence theory are especially valuable here, because the
macro-pattern cannot be reduced to any one participant's isolated intention.
\end{discussion}
